\answer{ex:10:1}{(а)~$\approx11.4$ бит на импульс, $\approx10^{12}$ бит/с. (б)~$\approx40$ бит/с, в $\sim10^{10}$ раз меньше. (в)~$2.5\cdot10^{10}\,\text{с}$ --- около восьмисот лет. (г)~$10$--$100$ мВт, $0.1$--$1\,\%$ энергии сигналов --- светодиод индикатора.}
\answer{ex:10:2}{(а)~(i)~$g\,w_{EI}w_{IE}>g(1-a)w_{EE}-1$; (ii)~$\tau_I<\tau_E/\bigl(g(1-a)w_{EE}-1\bigr)$. (б)~$g=4$, $\approx62$~Гц. (в)~$a>1/3$; при $a=0$ нарастающие колебания $\approx67$~Гц. (г)~Нет: устойчивость задаёт произведение $g(1-a)$.}
\answer{ex:10:3}{(а)~$\sigma^2R'_jx_j$. (б)~$x_j^2\sigma^4\bigl(\sum_{k\ne j}R'^2_k+2R'^2_j\bigr)$. (в)~$1/(N+1)$ на попытку, нужно $n\approx z^2(N+1)\propto N$ попыток. (г)~$0.51$, $0.197$, $0.059$, $0.016$ против $0.5$, $0.2$, $0.059$, $0.015$.}
\answer{ex:10:4}{(а)~$L_v\in[32.9,\,47.9]\ms$, $d=35.4\ms$, $\Delta_{\max}=7.5\ms$. (б)~$\tau_D=\Delta_{\max}/\ln2=10.8\ms$. (в)~$2\Delta_{\max}r_ar_v\approx0.38$ в секунду, пропорционально $r_ar_v$ и разбросу задержек.}
\answer{ex:10:5}{(а)~$\alpha=0.2$, ошибка $0.180$. (б)~$q=0.03$, $0.180$ --- то же. (в)~$\lambda=0.9$, $c=20$: $0.098$, на $45\,\%$ меньше. (г)~$0.025$; мешают запаздывание обнаружения и пропуск малых скачков.}
\answer{ex:10:6}{(а)~$0.2+0.6(1-\alpha)^n$. (б)~$n=\ln\bigl(\ln1.5/(0.6g)\bigr)/\ln(1-\alpha)$: $24$ и $4$. (в)~$11$ и $2$. (г)~$Q_A\to Q_B$ и $P(A)\to\frac12$; узнать о $B$ можно только пробой --- ручка \nt{НОРА}.}
\answer{ex:10:7}{(а)~$S'=S\bigl(1-4\eta\sigma^2+4\eta^2\sigma^4(G+2)\bigr)$, лучшее $\eta\sigma^2=1/(2(G+2))$, $n\approx(G+2)\ln(1/\varepsilon)$. (б)~При $N=128$: $595$, $151$, $12$ попыток против $599$, $157$, $14$. (в)~$\approx10^6$ попыток, около месяца. (г)~Нужно $K\propto N$ --- по сути свой сигнал ошибки у каждого нейрона.}
